\documentclass[11pt]{article}
\usepackage[a4paper,margin=18mm]{geometry}
\usepackage{fontspec}
\setmainfont{Latin Modern Roman}
\usepackage{amsmath,amssymb,amsthm,amscd,graphicx,color}
\usepackage{xurl}
\usepackage[unicode,hidelinks]{hyperref}
\allowdisplaybreaks
\setlength\emergencystretch{3em}
%\usepackage{amsmath}
%\usepackage{amsfonts}
%\usepackage{amssymb}
%\usepackage{amscd}
%\usepackage{graphics}
%\usepackage{epsfig}
%\input{epsf.sty}



%\textwidth5.6 in

%\newlength{\ghost}

%\raggedbottom

\newtheorem{thm}{Theorem}[section]
\newtheorem{thm*}{Theorem}
\newtheorem{lem}[thm]{Lemma}
\newtheorem{prop}[thm]{Proposition}
\newtheorem{cor}[thm]{Corollary}
\newtheorem{claim}[thm]{Claim}
\newtheorem{defprop}[thm]{Definition-Proposition}
\newtheorem{conj}[thm]{Conjecture}
\newtheorem{fact}[thm]{Fact}

\newtheorem{introthm}{Theorem}
\newtheorem{introprop}{Proposition}

\renewcommand{\theintrothm}{\Alph{introthm}}
\renewcommand{\theintroprop}{\Roman{introprop}}

{\theoremstyle{definition}

\newtheorem{df}[thm]{Definition}
\newtheorem{ex}[thm]{Example}
\newtheorem{qu}[thm]{Question}
\newtheorem{rmk}[thm]{Remark}
\newtheorem{nota}[thm]{Notation}
\newtheorem{assump}{Assumption}}

\def\nn{\nonumber}


\def\calB{{\mathcal B}}
\def\calP{{\mathcal P}}
\def\M{{\mathcal M}}
\def\Mco{{{}^{\rm c/o}\mathcal M}_{g,\delta_1,\dots,\delta_n}^{s,\beta}}
\def\Mcos{{{}^{\rm c/o}\mathcal M}_{g,\delta_1,\dots,\delta_n}^{s_1\dots s_k,\beta}}
\def\Sull{\rm {\mathcal S}ull^{\rm c/o}}
\def\wSull{\widetilde{\rm {\mathcal S}ull}^{\rm c/o}}
\def\Rp{{\mathbb R}_{>0}}
\def\la{\langle}
\def\ra{\rangle}

%\def\empty{\varnothing}

\def\idagger{r^{\dagger}}
\def\rdagger{r^{\dagger}}
\def\res{r}
\def\oto{\otimes \dots \otimes}
\newcommand{\aoa}[2]{a_{#1}\otimes \dots \otimes a_{#2}}
\newcommand{\bob}[2]{b_{#1}\otimes \dots \otimes b_{#2}}
\newcommand{\ata}[2]{a_{#1}\cdots a_{#2}}
\newcommand{\aca}[2]{a_{#1}, a_{#2}}
\newcommand{\set}[2]{\{#1,\dots, #2\}}

\def\wt{\omega}

\def\CO{{\mathcal O_\mu}}
\def\calO{{\mathcal O}}
\def\a{\alpha}
\def\t{\tau}
\def\B{\mathcal B}
\def\Z{\mathbb Z}
\def\R{\mathbb R}
\def\ba{\mathbf{a}}
\def\d{d^{{\rm CH}^*}}
\def\s{s^{{\rm CH}^*}}

%% Next two commands enlarge space for figures on the top of a page
\renewcommand{\textfraction}{0.01}
\renewcommand{\topfraction}{0.99}



\makeatletter
\def\tsep#1{%
  {\@tempdima=\ht\strutbox\advance\@tempdima #1
   \vrule height \@tempdima depth 0pt width 0pt\nobreak\hspace{0pt}%
  }%
}
\def\bsep#1{%
  {\@tempdima=\dp\strutbox \advance\@tempdima #1
   \nobreak\hspace{0pt}\vrule height 0pt depth \@tempdima width 0pt
  }%
}
\makeatother
\begin{document}
\begin{center}\Large Discrete brane-labelled open and closed arc-graph correlators respect all window gluings for basic ordinary Frobenius systems with commutative closed algebra and Euler compatibility\end{center}
\noindent\textbf{Stable ID:} 2757\par
\noindent\textbf{Authors:} Ralph M. Kaufmann\par
\noindent\textbf{Original work:} Open/Closed String Topology and Moduli Space Actions via Open/Closed Hochschild Actions\par
\noindent\textbf{Version:} Official SIGMA 6 (2010) article 036, DOI 10.3842/SIGMA.2010.036; journal PDF and native TeX acquired 2026-10-01, exact bytes pinned by SHA256\par
\noindent\textbf{Selected task scope:} Theorem5.1 restricted to the source basic brane-label system of ordinary finite-dimensional Frobenius algebras with algebra restriction maps, higher-intersection algebras zero, commutative A\_empty (C) and Euler compatibility (E), as defined4.1.4 and4.5. For discretely positively weighted brane-labelled arc graphs, original correlators(4)-(5) are compatible with all permitted matching-weight open and closed window gluings, including self and consecutive-window gluing, via the original Casimir contractions. No dg-PROP or moduli-space conclusion, general propagator systems or removal of C/E.\par
\noindent\textbf{Difficulty:} H2: Open window gluing can change the number of complementary regions and create labelled punctures, so the closed Frobenius contraction formula alone does not provide an action. The author defines weights for decorated boundary pieces and punctures and checks every local non-self and self-gluing topology. The new lone-window cases require the projection and Euler compatibility identities, with explicit contractions producing exactly the required Euler and puncture factors.\par
\noindent\textbf{Editorial boundary note (not author text):} The complete original Theorem5.1 is retained as the finite discrete arc-graph correlator compatibility outcome, with the basic brane-label Frobenius system, commutativity(C) and Euler compatibility(E) used in the author proof. The original window/brane labels, matching-weight gluing, boundary decoration, puncture/Euler weights, Casimir composition and every non-self, self and consecutive-window contraction are bundled; precise c/o operations are included from AppendixA. Ordinary Frobenius inputs and cited closed Hochschild sign conventions remain foundations. Later dg-PROP/moduli-space extensions and general propagator systems are unselected.\par
\noindent\textbf{Source:} \url{https://sigma-journal.com/2010/036/}\quad DOI: \url{https://doi.org/10.3842/SIGMA.2010.036}\par
\noindent\textbf{License and attribution:} Copyright retained by the named authors. Published by SIGMA under CC BY 4.0: \url{https://creativecommons.org/licenses/by/4.0/}. Source: \url{https://sigma-journal.com/about.html}. Adaptation consists of standalone excerpt formatting, cover and numbering restoration; original author language and mathematical text are retained.\par
\noindent\textbf{Editorial symbol notice (not author text):} The printed brief Theorem5.1 states basic Frobenius systems; the actual proof expressly uses commutativity(C) and Euler compatibility(E). The selected claim retains both authored hypotheses, and uses ordinary algebra inputs rather than separately counting a super-sign extension. All original formulas and wording remain unchanged.\par
\medskip\hrule\medskip\noindent\textbf{Author text begins below.}\par
\setcounter{section}{2}
\setcounter{subsection}{0}
\setcounter{subsubsection}{0}
\setcounter{thm}{0}
\setcounter{equation}{0}
% Original source lines 275--410
\subsection{Arc graphs in brane labelled windowed surfaces}

A {\it windowed surface} $F=F^s_g(\delta_1,\ldots ,\delta _r)$ is a
smooth oriented surface of genus $g\geq 0$ with $s\geq 0$ punctures and
$r\geq 1$ boundary components together with the specif\/ication of a
non-empty f\/inite subset $\delta _i$ of each boundary component,
for $i=1,\ldots ,r$, and we let $\delta =\delta
_1\cup\cdots\cup\delta _r$ denote the set of all distinguished points in the boundary
$\partial F$ of $F$ and let  $\sigma$ denote the set of all punctures.
The set of components of
$\partial F-\delta$ is called the set $W$ of {\it windows}.

Furthermore one needs to specify a brane labelling $\beta$.
For this we f\/irst f\/ix a set $\calB$ of basic brane labels
and denote by $\calP(\calB)$  power set.  Notice that $\varnothing\in\calP(\calB) $,
this will encode the closed sector. The elements of $\calP(\calB)$ of cardinality
bigger than one should be thought of as intersecting branes. These
give room for extra data, but it is possible to set all the contributions
for these  ``higher intersection'' branes
to zero in a given model.


A {\it brane-labeling} on a windowed surface $F$ is a function
\[
\beta : \ \delta\coprod\sigma\to{\mathcal P}({\mathcal B}),
\]
where $\sqcup$ denotes the disjoint union, so that if $\beta
(p)=\varnothing$ for some $p\in\delta$, then $p$ is the unique point
of $\delta$ in its component of~$\partial F$.  A brane-labeling
may take the value $\varnothing$ at a puncture.



A window $w\in W$ on a windowed surface $F$ brane-labeled by
$\beta$ is called {\it closed} if the endpoints of $w$ coincide at
the point $p\in\delta$ and $\beta (p)=\varnothing$; otherwise, the
window $w$ is called {\it open}.
Each window def\/ines a pair of brane labels $\beta(w)$ which is
the pair $(S,T)$ of the brane labels of the beginning and end of the window (these
may coincide).

\subsubsection{Arc families}

We def\/ine the sets
\[
\delta
(\beta )=\{ p\in\delta :\beta (p)\neq\varnothing\},\qquad \sigma
(\beta )=\{ p\in\sigma :\beta (p)\neq\varnothing\} .
\]


Def\/ine a $\beta$-arc $a$ in $F$ to be an arc properly embedded in $F$ with its endpoints in $W$ so that~$a$ is not
homotopic, f\/ixing its endpoints, to $\partial F-\delta (\beta )$.  For example, given a distinguished point
$p\in\partial
F$, consider the arc lying in a small neighborhood that simply connects one side of $p$ to another in $F$; $a$ is
a $\beta$-arc if and only if
$\beta (p)\neq\varnothing$. We will call such an arc a {\it small arc} around $p$.

Two $\beta$-arcs are {\it parallel} if they are homotopic rel~$\delta$, and a $\beta$-arc family is the homotopy
class
rel~$\delta$ of a collection of $\beta$-arcs, no two of which are parallel.  Notice that we take homotopies rel~$\delta$
rather than
rel $\delta (\beta )$.





Given a positively weighted arc family in $F$, let us furthermore say that a window $w\in W$ is {\it active}
if there
is an arc in the family with an endpoint in $w$, and otherwise the window is {\it inactive}.



\subsubsection{The mapping class group and arc graphs}

The {\it $($pure$)$ mapping class group} $MC(F)$ of $F$ is the group of
orientation-preserving homeomorphisms of $F$ pointwise f\/ixing
$\delta\cup\sigma$ modulo homotopies pointwise f\/ixing $\delta\cup\sigma$.


$MC(F)$ acts naturally on the set of $\beta$-arc families.
An {\it arc graph} is an equivalence graph under this action.

\subsection{Arc spaces, moduli spaces and the open/closed Sullivan PROP}
\subsubsection{The Arc spaces}



A {\it  weighting} on an arc family is the assignment of
a positive  real number to each of its components.
A weighting naturally passes to the arc graph.
A weighting is called {\it discrete} if it takes values in ${\mathbb N}$.



Let ${\rm Arc}'(F,\beta )$ denote the geometric realization of the partially
ordered set of all $\beta$-arc families in $F$.  ${\rm Arc}'(F,\beta )$ is
described as the set of all projective positively weighted $\beta$-arc
families in $F$ with the natural topology. (See for instance \cite{KLP} or \cite{P1} for
further detail.)

$MC(F)$ again
acts naturally on ${\rm Arc}'(F,\beta )$.  The {\it arc
complex} is def\/ined to be the  quotient under this action
\[
{\rm Arc}(F,\beta )={\rm Arc}'(F,\beta )/MC(F).
\]

We shall also consider the corresponding deprojectivized versions:
$\widetilde{\rm Arc}{}'(F,\beta )\approx {\rm Arc}'(F,\beta )\times{\mathbb R}_{>0}$ is the space of all positively weighted arc
families in $F$ with the natural topology, and
\[
\widetilde{\rm Arc}(F,\beta )=\widetilde{\rm Arc}{}'(F,\beta )/MC(F)\approx
{\rm Arc}(F,\beta )\times{\mathbb R}_{>0}.
\]



For any windowed surface $F$, def\/ine
\[
\widetilde{\rm Arc}(F)=\bigsqcup \widetilde{\rm Arc}(F,\beta),
\]
where the disjoint union is over all brane-labellings on $F$.
\[
\widetilde{\rm Arc}(n,m)=\bigsqcup \left\{ \alpha\in \widetilde{\rm Arc}(F):
\begin{array}{l}\alpha ~{\rm has}~n~{\rm closed~and}~ m~{\rm
open~active}\\{\rm
windows~and~no~inactive~windows}\end{array}\right\},
\] where the
disjoint union is over all orientation-preserving
homeomorphism classes of windowed surfaces.




\setcounter{section}{2}
\setcounter{subsection}{4}
\setcounter{subsubsection}{1}
\setcounter{thm}{0}
\setcounter{equation}{0}
% Original source lines 475--573
\section{Geometric c/o structures}
\label{geosection}

In \cite{KP} we axiomatized a structure of spaces that allow four dif\/ferent types of gluing: i.e.\
either open or closed and either self or non-self gluings together with a coloring of the open gluings by brane labels~$\beta$. In the above case the gluing is on two  windows, which are either both
open or closed and either belong to the same (self gluing) or disjoint surfaces (non-self gluing).
The coloring is by the brane labels. The color of a window is the pair of brane labels of its
end points. When we glue $\beta$ arc families, we glue the surfaces and
glue the arc families as foliations. We will now review this process according to~\cite{KP}.


\subsection{The gluing underlying the topological c/o structure}




If $\alpha\in\widetilde{\rm Arc} (n,m)$, then def\/ine the {\it $\alpha$-weight} $\alpha(w)$ of an active window $w$ to be the sum
of
the weights of arcs in $\alpha$ with endpoints in $w$, where we count with multiplicity
(so if an arc in $\alpha$ has both endpoints in $w$, then the weight of this arc contributes twice to the weight of
$w$).

Suppose we have a pair of arc families $\alpha_1$, $\alpha _2$
in respective windowed surfaces $F_1$, $F_2$ and a~pair of active windows $w_1$ in $F_1$ and $w_2$ in $F_2$, so that the
$\alpha _1$-weight of $w_1$ agrees with the $\alpha _2$-weight of $w_2$.
Since $F_1$, $F_2$ are oriented surfaces, so are the windows $w_1$, $w_2$ oriented.   In each operation, we
identify windows reversing orientation, and we identify certain distinguished points.

To def\/ine the open and closed gluing ($F_1\neq F_2$) and self-gluing ($F_1=F_2$) of $\alpha_1$, $\alpha _2$ along the
windows
$w_1$, $w_2$, we
identify windows and distinguished points in the natural way and combine foliations.

A crucial dif\/ference  between the closed and open string operations is that in the closed case,
the points are thought of as marked, which in the open case the points behave like punctures.
This means that in the closed case, we
replace the  distinguished point by simply forgetting that is was distinguished. This way  no puncture is created. In the open case the distinguished
points
always give rise to other distinguished points or perhaps punctures.
 In any case whenever distinguished points are
identif\/ied, one takes the union of brane labels (the intersection of branes) at the new resulting distinguished point or puncture.

\begin{figure}[th]
%\epsfxsize = \textwidth \epsfbox{stropfig4.eps}
\centerline{\includegraphics{Kaufmann-Fig1}}

\caption{The dif\/ferent cases of gluing.}\label{stropfig4}
\end{figure}


\subsubsection{Closed gluing and self-gluing}
Identify the two corresponding boundary components of $F_1$ and
$F_2$, identifying also the distinguished points on them {\it and
then including this point in the resulting surface} $F_3$. $F_3$~inherits a~brane-labeling from those on $F_1$, $F_2$ in the natural
way.  We furthermore glue $\alpha _1$ and $\alpha _2$ together in
the natural way, where the two collections of foliated rectangles
in~$F_1$ and~$F_2$ which meet $w_1$ and $w_2$ have the same total
width by hypothesis and therefore glue together naturally to
provide a measured foliation ${\mathcal F}$ of a closed subsurface
of $F_3$.


\subsubsection{Open gluing}

The surfaces $F_1$ and $F_2$ are distinct, and we identify $w_1$ to $w_2$ to produce
$F_3$.
There are cases depending upon whether the closure of $w_1$ and $w_2$ is an interval or a circle.  The salient
cases are
illustrated in Fig.~\ref{stropfig4},~b--d.  In each case,  distinguished points on the boundary in~$F_1$ and~$F_2$ are
identif\/ied to produce a new distinguished boundary point in $F_3$, and the brane labels are combined, as is also
illustrated.  As before, since the
$\alpha _1$-weight on $w_1$ agrees with the $\alpha _2$-weight on $w_2$, the foliated rectangles again combine to
provide a measured foliation~${\mathcal F}$ of a closed subsurface of~$F_3$.

 {\bf Open self-gluing.}~There are again cases depending upon whether the closure of $w_1$ or $w_2$ is a~circle or an interval, but there is a further case as well when the two intervals lie in a common boundary component
and
are consecutive.  Other than this last case, the construction is identical to those illustrated
in Fig.~\ref{stropfig4},~b--d.  In case the two windows are consecutive along a common boundary component, again they are identif\/ied
so
as to produce a surface~$F_3$ with a puncture resulting from their common endpoint as in Fig.~\ref{stropfig4},~e--f, where the
puncture is
brane-labeled by the label of this point, and the foliated rectangles combine to provide a measured foliation~${\mathcal F}$ of a closed subsurface of~$F_3$.



At this stage, we have only constructed a measured foliation ${\mathcal F}$ of a closed subsurface of $F_3$, and
indeed,
${\mathcal F}$ will typically not be a weighted arc family, but the sub-foliation
${\mathcal
F}'$ comprised of leaves that meet $\partial F$ corresponds to a weighted arc family $\alpha _3$ in $F_3$.
Notice that
the $\alpha _3$-weight of any window uninvolved in the operation agrees with its $\alpha _1$- or $\alpha _2$-weight.

The assignment of $\alpha_3$ in $F_3$ to $\alpha _i$ in $F_i$, for
$i=1,2$ completes the def\/inition of the various operations.
Associativity and equivariance for bijections are immediate, and so
we have our f\/irst non-trivial example of a c/o structure; see
Appendix~\ref{appendixA} for the precise def\/inition.

\setcounter{section}{3}
\setcounter{subsection}{2}
\setcounter{subsubsection}{0}
\setcounter{thm}{3}
\setcounter{equation}{0}
% Original source lines 635--1412
\subsection[Discretization of the model ${\mathbb N}$-valued foliations]{Discretization of the model $\boldsymbol{{\mathbb N}}$-valued foliations}

We wish to point out that the subset of discrete valued foliations is stable under the compositions
in both the c/o structure and the c/o PROP structure.


\subsubsection[Discrete representation for discretely weighted $\beta$-arc families]{Discrete representation for {\em discretely weighted} $\boldsymbol{\beta}$-arc families}

In order to decorate, we will change the picture slightly. Previously we had arc graphs, whose
edges are not allowed to be parallel. For a {\em discretely weighted} $\beta$-arc families with weighting $wt$
we  will consider its {\em leaf representation} to be the foliation each of whose bands $e$ has
$wt(e)$ number of leaves. This means that we consider an new type of arc graph which has $wt(e)$ parallel
edges for each underlying edge of the original arc graph. We will call this the {\em discrete representative}
of the arc graph.


\section{Algebraic c/o structures}
\label{algsection}

Given a $\beta$-arc graph with weights that are natural numbers, we are going to
associate an operation on the Hochschild complexes  ${\rm CH}^*(A_{\varnothing},M_{B,B'})$
of a f\/ixed Frobenius algebra $A_{\varnothing}$ with coef\/f\/icients in a module $M_{B,B'}$ with $B,B'\in \calB$.
These modules will be given by tensor products of Frobenius algebras $A_B$ indexed by elements of $\calB$: $M_{B,B'}: A_{B'}\otimes A_B$. These operations will be def\/ined on the isomorphic
double sided bar-complexes  $B(A_B,A_{\varnothing},A_{B'})$. For most operations we will need
that $A_{\varnothing}$ is commutative, but this is not always the case.  We will indicate
when the commutativity can be dropped.


\subsection{Frobenius algebras and systems of Frobenius algebras}

\subsubsection{Notation for Frobenius algebras}
The main actors are  Frobenius algebras, so we will f\/ix some notation.

Recall that a Frobenius algebra (FA) is a triple $(A,1, \langle \; ,\;\rangle)$
where $(A,1)$ is a unital (super-) algebra and $\langle \; ,\;\rangle$ is a non-degenerate
(super-) symmetric even pairing which satisf\/ies
\begin{equation*}
\langle ab,c\rangle=\langle a,bc\rangle.
\end{equation*}

We will set
\begin{equation*}
\int a :=\langle a , 1 \rangle.
\end{equation*}
Then $\int$ is cyclically (super-)invariant, i.e.\ a trace
\begin{equation*}
\int abc = \la ab,c\ra =\la c, ab\ra=\int cab.
\end{equation*}


Since $ \langle \; ,\;\rangle$ is non-degenerate on $A$, so is $
\langle \; ,\;\rangle_{A\otimes A}:=\langle \; ,\;\rangle\otimes
\langle \; ,\;\rangle\circ \tau_{2,3}$ on $A^{\otimes 2}\otimes
A^{\otimes 2}$. Here~$\tau_{2,3}$ is the commutativity constraint
for the symmetric monoidal category applied to the second and third
factors, be it the category of vector spaces, dg-vector spaces or
$\Z/2\Z$ graded vector spaces
 In our current setup this just interchanges the second and third  factors of
$A\otimes A\otimes A\otimes A$  or in the super case changes these factors and introduces
the usual super sign.

We will omit all super signs from our discussion as they can be added in a straightforward fashion.

The multiplication $\mu:A\otimes A\to A$ has an adjoint $\Delta: A\to A\times A$
def\/ined by
\begin{equation}
\label{inveq} \langle \Delta a, b\otimes c\rangle_{A\otimes
A}=\langle a,bc\rangle.
\end{equation}

Moreover given a FA we will consider a basis $\Delta_i$, set $g_{ij}=\langle \Delta_i,\Delta_j\rangle$
and let $g^{ij}$ the coef\/f\/icients of the inverse of $(g_{ij})$, viz.\ the inverse ``metric''.

There are two special elements
\begin{equation*}
e=\mu\Delta(1)=\sum_{ij} \Delta_i g^{ij}\Delta_j,
\end{equation*}
which we call the Euler element and
\begin{equation*}
C=\Delta(1)=\sum \Delta_i g^{ij} \otimes \Delta_j,
\end{equation*}
which we call the Casimir element.

Notice that Euler element commutes with every element.
\begin{equation*}
ae=a\mu\Delta(1)=\mu\Delta(a)=\mu\Delta(1)a=ea.
\end{equation*}
This follows by direct computation and the Frobenius relations
\begin{equation*}
({\rm id}\otimes \mu)(\Delta\otimes {\rm id})=\Delta\mu=(\mu\otimes
{\rm id})({\rm id}\otimes \Delta),
\end{equation*}
which in turn follow from the def\/inition of $\Delta$ and the
invariance of the pairing~(\ref{inveq}).

\subsubsection{Adjoint maps}
Notice that for any map $r:A\to B$ between two Frobenius algebras
there is an {\em adjoint map} $r^{\dagger}:B\to A$ def\/ined by
\begin{equation*}
 \langle r^{\dagger}(b) ,a  \rangle =\langle b,r(a)\rangle.
\end{equation*}

This is equivalent to
\begin{equation}
\label{adjointeq}
 \int r^{\dagger}(b)a  =\int br(a).
\end{equation}
%In this sense $\Delta=\mu^{\dagger}$.

These maps arise
in geometric situations as follows. Let $i:N \to M $ be the inclusion map, where $M$ is a compact manifold and $N$ is a compact submanifold. Then $i$ induces a map $i^*$ from
$A:=H^*(M)$ to $B:=H^*(N)$. By Poincar\'e duality there is a push forward $i_*:B\to A$. In the previous notation if  $r=i^*$ then $r^{\dagger}=i_*$.


\begin{lem}In general, we have the Projection Formula
\begin{equation*}
r^{\dagger}(r(a)b)=ar^{\dagger}(b). % \qquad \mbox{\rm Projection Formula}
\end{equation*}
\end{lem}
\begin{proof}
\begin{gather*}
\la r^{\dagger}(r(a)b), c\ra = \la r(a)b,r(c)\ra = \la b, r(ca)\ra = \la ar^{\dagger}(b),c\ra,
\end{gather*}
where we used the cyclic symmetry of the product twice.
\end{proof}

With  the self-intersection condition Section~\ref{selfdef} in mind, we def\/ine the element
\begin{equation*}
e_r^{\perp}:=r (r^{\dagger}(1)).
\end{equation*}

\subsubsection{Systems of Frobenius algebras}
\begin{df} A $\calB$-Frobenius algebra is a set of Frobenius
algebras $A_S$ indexed by $S\in \calP(\calB)$ together with algebra maps
$r_{S,S'}:A_S\to A_S'$ whenever $S\subset S'$,
such that for $S\subset S'\subset S''$: $r_{S',S''}\circ r_{S,S'}=r_{S,S''}$.
\end{df}



Note that in particular if $A_{\varnothing}$ is commutative
every $A_S$ is an $A_{\varnothing}$ module via
the restriction map. More precisely,  every $A_S$ is a left and a right $A_{\varnothing}$
module via the maps $\lambda (a,a'):= r_{\varnothing S}(a)a'$
and $\rho(a,a') :=a' r_{\varnothing S}(a)$. If $A_{\varnothing}$ is not commutative,
we still have that $A_S$ is a left $A_{\varnothing}$ module and a right $A_{\varnothing}^{\rm op}$ module.


\subsubsection{Basic brane label systems}
Given a  brane label set $\calB$ one set choices of $\calB$-Frobenius algebras is given by
a collection~$A_B$, $B\in \calB$ and $A_{\varnothing}$ together with maps $r_B:A_{\varnothing}\to A_B$.

For any $S\in \calP\calB$ with $|S|\geq 2$ we simply set $A_S=0$ where
we allow the zero algebra to be a~Frobenius algebra.

We call these systems basic brane label systems and for simplicity deal only with these.
The data of the Frobenius algebras and morphisms will be called a basic $\calB$ Frobenius algebra.

\begin{rmk}
In the following we will mostly deal with only basic $\calB$ Frobenius algebra in order to not unduly
burden the  reader with yet more  structures that we would need in order to deal with the general case.
All the results do however generalize to the general case if we introduce propagators to f\/ix the
c/o structure on the algebraic side.
\end{rmk}

\begin{nota}
Given $A_B$ we will use the notation
$1_B$, $e_B$, $\Delta^i_B$, $\la$, $\ra_B$, $g^B_{ij}$, $g^{ij}_B$, $e^{\perp}_B$,
$\dots$ for its unit, Euler element,
basis $A_B$, metric inverse metric, $e^{\perp}_{r_B}$ etc.

We will also omit the label ${\varnothing}$ i.e.\ write  $e$ for $e_{\varnothing}$, $A$ for $A_{\varnothing}$ if no
confusion can arise.
\end{nota}



\begin{df}
\label{Eulerdef}
We say that a basic $\calB$-FA  satisf\/ies the condition of commutativity $(C)$ if $A_{\varnothing}$ is commutative.

And we say that a $\calB$-FA satisf\/ies the
the Euler compatibility condition or the {\em condition $(E)$} if for all $B\in \calB$, $a^{(1)},a^{(2)}\in A_B$
\begin{equation*}
(E) \quad \sum_{ij}\res_B^{\dagger}\big(a^{(1)}\Delta^B_i\big) g^{ij}_B
 \,  \res_B^{\dagger}\big(\Delta^B_ia^{(2)}\big)=e_{\varnothing}\res_B^{\dagger}\big(a^{(1)}a^{(2)}\big).
\end{equation*}
\end{df}


\begin{df}
\label{selfdef}
A basic $\calB$-Frobenius algebra satisf\/ies  the {\em self-intersection} condition $(I)$
if for all $B\in \calB$
\begin{gather*}
(I_1) \quad r_{B}r^{\dagger}_{B}  (a)=  a e^{\perp}_{B} \qquad \text{and}\qquad
(I_2) \quad  e_Be^{\perp}_B =  \res_B(e).
\end{gather*}
\end{df}



\begin{prop}A  basic system of  $\calB$ Frobenius algebra which satisfies the self-intersection condition $(I)$ satisfies the Euler condition $(E)$.
\end{prop}

\begin{proof} We will show that the r.h.s.\ and the l.h.s.\ of $(E)$ have
the same inner product with any element $b$ of $A$:

For all $b\in A_{\varnothing}$, $a, a'\in A_B$
\begin{gather*} \sum_{ij} \int b  \res_B^{\dagger}\big(a\Delta^{B}_i\big)
g_{B}^{ij}  \res_B^{\dagger}\big(\Delta^{B}_j a'\big)
 = \sum_{ij}g_{B}^{ij}
 \int_B \res_{B}(b)a\Delta^{B}_i
 \res_{B}\res_B^{\dagger}\big(\Delta^{B}_ja'\big) \\
\qquad{}  \stackrel{(I_1)}{=} \sum_{ij}g_{B}^{ij}
 \int_B \res_{B}(b)a\Delta^{B}_i \Delta^{B}_j
 a'e_{B}^{\perp}
 =
 \int_B \res_{B}(b)aa' e_B e_B^{\perp}\\
\qquad{} \stackrel{(I_2)}{=} \int_B \res_{B}(b)aa'\res_B(e)
 = \int b  e\res_B^{\dagger}(aa'),
\end{gather*}
where the f\/irst   equality follows from equation (\ref{adjointeq}),
the third equality from the def\/inition of $e_B$ and the fact that~$e_B$ as the Euler element commutes with all other elements and the
last equality follows from the projection formula and the fact that~$e$ commutes.
\end{proof}



\subsubsection{Geometric data}
\label{geodatasection} One example of the basic data is given by a
compact manifold $M$ together with a collection $N_B\subset M$, $B\in
\calB$ of compact submanifolds. We can then set $A_B:=H^*(N_B)$ and
use the restriction maps $r_B$     given by pullback. These satisfy
the f\/irst  condition $(I_1)$ of $(I)$ due to the self-intersection
formula where $e_B^{\perp}=e(N_{M/N_B})$ is the Euler class of the
normal bundle of $N_B$ in $M$. The second condition ($I_2$) follows
from the excess intersection formula for homology~\cite{Qu} applied
to the diagram
\begin{equation}
\label{intdiagram}
\begin{CD}
N_B @>\Delta_B>>N_B \times N_B \\
@Vi_BVV @V i_B\times i_B VV\\
M@>\Delta>> M\times M
\end{CD}
\end{equation}
keeping in mind that $\mu=\Delta^{*}$, $\res_B=i_B^*$
\begin{gather*}
\res_B(e)=\res_B\mu\mu^{\dagger}(1)=i_B^*\Delta^*\Delta_*(1)=\Delta^*_B(i_{B}^*\times i_{B}^*)\Delta_*(1) \\
\phantom{\res_B(e)}{} =e_B^{\perp}\Delta_B^*\Delta_{B*}i_B^*(1)=e_B^{\perp}\Delta_B^*\Delta_{B*}(1)=
e^{\perp}_Be_B.
\end{gather*}

Alternatively one can use decomposition $TM|_{N_B}=TN_B\oplus N_{M/NB}$ and multiplicativity of the Euler class.

\begin{rmk}
We can also formalize the geometricity by  staying in the framework of Frobenius algebras, but
postulating a new axiom which guarantees the equations one would obtain from excess intersection formulas from all embedding and intersection diagrams analogous to~(\ref{intdiagram}).
\end{rmk}

\begin{cor}
A basic geometric $\calB$ Frobenius algebra satisfies the Euler condition $(E)$.
\end{cor}

\subsubsection{Hochschild complexes}
The action on the closed sector is on the Hochschild cochain-complex of $A$.
Recall that the Hochschild chain complex of an $A$ bimodule $M$ is the complex ${\rm CH}_n(A,M)$
\[
{\rm CH}_n(A,M)=M\otimes A^{\otimes n}
\]
and whose  dif\/ferential is given by $d=\sum_i (-1)^i d_i$,
where the $d_i$ are the pre-simplicial dif\/ferentials
\begin{gather*}
d_0 (m\otimes a_1 \oto a_n ) = m a_1 \otimes\aoa{2}{n},\\
d_i (m\otimes a_1 \oto a_n) = m\otimes \aoa{1}{i-1}\otimes a_{i}a_{i+1}\otimes
\aoa{i+2}{n}, \\
d_{n} ( m\otimes a_1 \oto a_n ) = a_n m\otimes\aoa{1}{n-1}.
\end{gather*}


There are degeneracies inserting $1$ into the $i$th position
\begin{gather*}
s_i: \ {\rm CH}_n(M,A) \to {\rm CH}_{n+1}(M,A),\\
\phantom{s_i:}{} \ \ m\otimes a_1 \oto a_n \mapsto m\otimes a_1
\oto a_{i-1}\otimes  1 \otimes a_i \oto a_n.
\end{gather*}

The Hochschild chain complex ${\rm CH}_*(A,M)$ is also sometimes called
the cyclic bar complex and is denoted by $B_*(A,A)$.


The Hochschild co-chain complex is dually given by
\[
{\rm CH}^n(A,M)={\rm Hom}(A^{\otimes n},M)
\]
with the dual dif\/ferential $\d=\sum_i (-1)^i \d_i$,
$\d:{\rm CH}^n(A,M)\to {\rm CH}^{n+1}(A,M)$, where for $f\in {\rm CH}^n(A,M)$
\begin{gather*}
\d_0 f( a_1 \oto a_{n+1} ) =  a_1f(\aoa{2}{n}),\\
\d_i f(a_1 \oto a_{n+1}) = f(\aoa{1}{i-1}\otimes a_{i}a_{i+1}\otimes
\aoa{i+2}{n}), \\
\d_{n+1} ( a_1 \oto a_{n+1}) = f(\aoa{1}{n})a_{n+1}.
\end{gather*}

The degeneracies dualize to $\s_i:{\rm CH}^n(A,M)\to {\rm CH}^{n-1}(A,M)$
\begin{equation*}
\s_i f(a_1 \oto a_{n-1})=f(\aoa{1}{i-1}\otimes 1\otimes
\aoa{i}{n-1}).
\end{equation*}


In case $M=A$ multiplication of functions gives a natural product
\[
\cup: \ {\rm CH}^n(A,A)\otimes {\rm CH}^m(A,A)\to {\rm CH}^{n+m}(A,A).
\]


Notice that if $A$ is a Frobenius algebra, it is isomorphic to its
dual as a bi-module. Since $A\simeq {\check A}$, ${\rm CH}^n(A,A)\simeq
A\otimes {\check A}^{\otimes n}\simeq A^{\otimes n+1}\simeq
{\rm CH}_n(A,A)$. The cup product  and the product pairing, make
${\rm CH}^{\bullet}$ into a graded Frobenius algebra. Furthermore, the
dif\/ferentials  $d$ and $\d$ dualize to one and another.

\subsubsection{Reduced Hochschild complex}

For technical reasons discussed in \cite{hoch2} it is actually
easier to work with the reduced Hochschild co-chain complex
$\overline{\rm CH}^*(A,A)$. This complex is the subcomplex of functions
that vanish on all degeneracies that is functions $f:A^{\otimes
n}\to A$ such that $f(a_1,\dots,1,\dots, a_n)=0$ where $1$ is
plugged in into any position. The complex inherits the dif\/ferential
and computes the same cohohmology as the original complex. Dually
there is the reduced bar complex $\overline {\rm CH}_*(A,A)$ or $\bar
B_*(A,A)$ which we shall use in the closed sector.


\subsubsection{Double sided bar construction}
For the open action,
we will consider the double sided bar complexes
$B(S,T):=B_{\bullet}(A_S,A,A_T)$
whose components are def\/ined as
\[
B_n(A_S,A,A_T)=A_S\otimes A^{\otimes n}\otimes A_T
\]
and whose dif\/ferential is given by $d=\sum_i (-1)^i d_i$, where the
$d_i$ are the pre-simplicial dif\/ferentials
\begin{gather*}
d_0 (a_S\otimes a_1 \oto a_n \otimes a_T) = a_Sa_1\otimes \aoa{2}{n}\otimes a_T,\\
d_i (a_S\otimes a_1 \oto a_n \otimes a_T) = a_S\otimes
\aoa{1}{i-1}\otimes a_{i}a_{i+1}\otimes
\aoa{i+2}{n} \otimes a_T,\\
d_{n} (a_S\otimes a_1 \oto a_n \otimes
a_T) = a_S\otimes\aoa{1}{n}\otimes a_n a_T.
\end{gather*}

\begin{rmk}
Notice that since we are dealing with Frobenius algebras, this is
isomorphic to  ${\rm CH}^{\bullet}(A,A_T\otimes A_S)$ and again the
dif\/ferentials dualize to one and another.
\end{rmk}

\subsubsection{Degeneracies}
The double sided complex is actually simplicial,
which means that it also has degeneracy maps
\begin{gather*}
s_i: \ B_n(A_S,A,A_T) \to B_{n+1}(A_S,A,A_T), \\
\phantom{s_i:} \ \ a_S\otimes a_1 \oto a_n \otimes a_T \mapsto a_S\otimes a_1 \oto
a_{i-1}\otimes  1 \otimes a_i \oto a_n \otimes a_T.
\end{gather*}

\subsubsection{Brane labelled  bar complexes}

Fix a $\calB$ Frobenius algebra. For a window $w$ on a brane
labelled surface with labelling $\beta$ we set $B(w):=B(\beta (w))$
if $w$ is open and
$B(w)=B(\beta(w))=B(\varnothing,\varnothing):=\overline{\rm CH}_n(A,A)$ if $w$ is
closed.

\subsection{Gluing in brane labelled complexes}

As we have noted, the complexes $B(\beta)$ have non-degenerate
graded inner products which allows us to dualize them.
 Using this
dualization, we can compose two correlation functions for the basic
brane label case. In the general case we would need to introduce
propagators to do this. We will refrain from adding this technical
point here for clarity of the discussion.

Given  (graded) vector spaces $V_{\beta}: \beta\in I$ over a ground
f\/ield $k$ with an involution~$\bar{}$ on~$I$, isomorphisms
$V_{\beta}\simeq V_{\bar\beta}$ s.t.\ $\bar{\bar{a}}=a$ and (graded)
non-degenerate even symmetric pairings $\la\;,\;\ra_{\beta}$ for
each $V_{\beta}$
 let $C_{\beta}=\sum
\Delta^{\beta}_i g^{ij}_{\beta} \otimes \bar\Delta^{\beta}_j$ be the
Casimir element for the induced pairing between $V_{\beta}$ and
$V_{\bar\beta}$ expressed in a basis $(\Delta^{\beta}_i)$.


We can compose two correlators two correlators $Y:\bigotimes_{l \in
L} V_{\beta_l} \to k$ and $Y':\bigotimes_{l'\in L'}
V_{\beta_{l'}}\to k $ where $L$ and $L'$ are labelling sets by
inserting a Casimir.

More precisely given $l_0\in L$ and $l'_0\in L'$ such that
$\beta_{l_0}=\bar\beta_{l_0'}=\beta$ we def\/ine their composition
$Y\circ_{l_0l_0'} Y': \bigotimes_{l''\in (L \setminus \{l_0\} \sqcup
L' \setminus \{l_0'\})} \to k$ by
\begin{gather*}
Y\circ_{l_0l_0'} Y' \Big(\bigotimes_{l\in L\setminus \{l_0\}} v_l\otimes
\bigotimes_{l'\in L'\setminus \{l_0'\}} v_{l'}\Big) \\
\qquad{}
=\sum_{ij}
Y\bigg(\Big(\bigotimes_{l\in L\setminus \{l_0\}} v_l\Big)\otimes
\Delta^{\beta}_i\bigg)g^{ij}Y'\bigg(\Big(\bigotimes_{l'\in L'\setminus
\{l_0'\}} v_{l'}\Big)\otimes \bar \Delta^{\beta}_j\bigg),
\end{gather*}
where $\Delta^{\beta}_i$ is in the position $l_0$ and $\bar
\Delta^{\beta}_j$ is in position $l_0'$. If $\beta_{l_0}\neq
\bar\beta_{l'_0}$ we set the composition to $0$. In the more general
case these would be non-zero and the composition would use
propagators.


Alternatively, we could also dualize the maps $Y$ using
$C_{\beta_l}$ in any position $l$ to obtain maps to~$V_{\beta_l}$
instead of $k$ and then compose these maps. There is an analogous
procedure for self-gluing.

\begin{rmk}
In our case $I=\calB\times \calB\cup \{\varnothing,\varnothing\}$,
$\overline{(S,T)}=(T,S)$, $V_{\beta}=B(\beta)$ the corresponding bar
complex and the isomorphisms $\bar{}:B(S,T)\to B(T,S)$  are given by
\begin{equation*}
a_S\otimes \aoa{1}{n}\otimes a_T\mapsto a_T \otimes
\aoa{n}{1}\otimes a_S
\end{equation*}
and $B(\varnothing,\varnothing)\to B(\varnothing,\varnothing)$
\begin{equation*}
 \aoa{0}{n} \mapsto a_0 \otimes
\aoa{n}{1}.
\end{equation*}
\end{rmk}





\begin{rmk}
In order to give a brane labelled c/o structure, we should consider
slightly enriched more complicated data. For this we would have to
look at tensors products of cyclic tensor products of bar complexes
$B(S_1,S_2)\otimes B(S_2,S_3)\oto B(S_n,S_1)$. Again in the interest
of brevity, we will not introduce this kind of complexity in a
formal fashion here.
\end{rmk}

\begin{rmk}
There are actually two c/o structures, one can compose the bar
complexes or their duals, viz.\ the correlators. Of course these
operations are dual to each other.
\end{rmk}


\section{Correlators}
\label{corsection}

\subsection{A universal formula for  correlators}
 There is a universal formula for the correlators. It is given by partitioning, decorating
and  de\-com\-posing the surface of a discretely weighted arc family
  along the arcs into little pieces of surface $S_i$ and integrating around
these pieces. We will now give the details.



\subsubsection{Decorating the boundary}
Fix a basic $\calB$ collection of Frobenius algebras.  For each {\em
discretely weighted} $\beta$-arc family $\alpha$,  we will def\/ine a
map
\[
Y(\alpha):= \bigotimes_{w \in \mathrm {Windows~of~}\alpha}B(\beta(w))\to k.
\]
These functions are homogeneous and their homogeneous components  are  zero by def\/inition unless
 $a_w\in B_{\alpha(w)-1}(\beta(w)) $.

Given a collection of homogeneous elements
we will decorate the pieces belonging to the boundary of the discrete representation of $\alpha$ by
the elements of the bar complexes
\begin{gather*}
a_w = a^w_S\otimes a_1^w \oto a^w_{\alpha(w)-1}\otimes a_T^{\prime w} \qquad \text{if} \quad \beta(w)=(S,T),\\
a_w = a^w_0\otimes a_1^w \oto a^w_{\alpha(w)-1}  \qquad \text{if} \quad
\beta(w)=(\varnothing,\varnothing).
\end{gather*}

Notice that  the boundary of the underlying surface minus the
discrete representative of the graph is a disjoint union of
intervals, which may or may not contain marked points. We call these
the {\em boundary pieces}. There are three types of boundary pieces
\begin{enumerate}\itemsep=0pt
\item[(1)] those not containing a marked point,
\item[(2)]  those
containing a marked point $\beta$-labelled by $\varnothing$,
\item[(3)] those containing a marked point with
$\beta$-label not $\varnothing$.
\end{enumerate}
In case~(3), if we remove the marked point we will have two components
which we will call half sides of the boundary piece. Now
each piece of type~(1) and the  half sides of the pieces of type~(3)
 belong to a unique window.
A piece of type~(2) comes from a unique closed window/boundary component. Moreover these
pieces all come in a natural linear order in each window as do the half sides of a piece of type (3) if
we consider the marked point to lie in between the half sides.


 We decorate
the boundary pieces as follows:

\vspace{1mm}

{\centering \begin{tabular}{@{\extracolsep{0pt}}c|l|l}
Type&description of $s$ &Decoration\bsep{1pt}\\
\hline
 (1) &$s$ is the $i$th piece of type (1) of the window $w$&$a_i^w\in A$\tsep{1pt}\bsep{1pt} \\
(2)& $s$ is the unique piece of type (2) of the closed window $w$&$a_0^w$ \bsep{1pt} \\
 (3)&the marked point
of  $s$ is labelled by $S$ and the half &
$(a_S^{\prime w_1},a_S^{w_2})$\\
&sides in their order belong to the &\\
&not necessarily distinct
 windows $w_1$, $w_2$&
\end{tabular}

}


%maybe A_empty=A or A_emtpy-C

\subsubsection{Weights}

 Let $\{S_i:i\in I\}$ be the components of
complement of the discrete representative of a given discretely weighted
$\beta$ arc family $\alpha$.
Each of these pieces has a polygonal boundary, where the sides of the polygons alternate between pieces
of the boundary and arcs running
between them. If we decorate the surface as described above every second side is a decorated
piece of boundary. We will call these the decorated sides.

To each decorated side $s$ of $S_i$, we associate a weight depending $\wt$
on its type and decoration.

\vspace{-2mm}

\begin{table}[h]\centering
\caption{General weights.}
\vspace{1mm}
\begin{tabular}{lll}
Type&Decoration& Weight $\wt(s)$\bsep{1pt}\\
\hline
(1) $s$ without marked point&$a\in A_{\varnothing}$&$a$\tsep{1pt}\bsep{1pt}\\
(2) $s$ with marked point marked by $\varnothing$&$a\in A_{\varnothing}$&$a$\bsep{1pt}\\
(3) $s$ marked point marked by $S$&$(a_S^{(1)},a_S^{(2)})$, $a_S^i\in A_{S}$&
$r^{\dagger}_{\varnothing S}(a_S^{(1)} a_S^{(2)})$
\end{tabular}
\end{table}





\subsubsection{The formula}
Moreover, the $S_i$ are oriented and so hence are their boundaries.
This means that the sides comprising each boundary component come with a cyclic order.
If there is only one boundary component for a given $S_i$, this gives a cyclic order over which we will integrate the given weights.
If there are more components, in which case the underlying arc family is not quasi-f\/illing,
then we need to assume (C) in order to make the following expression independent of choices.
For a homogeneous ${\mathbf  a}= \bigotimes_{w \in \mathrm {Windows~of~}\alpha} a_w \in \bigotimes_{w \in \mathrm {Windows~of~}\alpha}B(\beta(w))$ such that  $a_w\in B_{\alpha(w)-1}(\beta(w))$,
we decorate as above and def\/ine
\begin{equation}
\label{cordef}
Y_{S_i}({\mathbf a})=\int e^{-\chi(S)+1}\prod_{\substack{\text{Decorated sides}\\
\text{$s$ of $S_i$}}} \wt(s)  \prod_{\substack{\text{Punctures $p$}\\ \text{inside $S_i$}}}
r^{\dagger}_{\varnothing \beta(p)}(e_{\beta(p)}).
\end{equation}
If $\mathbf a$ is as above but there is  some $a_w\notin B_{\alpha(w)-1}(\beta(w))$ we set $Y_{S_i}(\mathbf a)=0$.

We then def\/ine{\samepage
\begin{equation}
\label{cordefgraph}
Y_{(\Gamma,w))}({\mathbf a}):=\prod_i Y_{S_i}({\mathbf a})
\end{equation}
and extend by linearity.}


\subsubsection{Signs} The correlators above actually have hidden signs which come
from the permutation of the input variables to their respective position. In the bar complexes these signs can be read of\/f by imposing that the tensor symbols have degree~1. On the geometric side
there are signs as well which are f\/ixed by f\/ixing an enumeration of the f\/lags, angles or edges.
In general we adhere to the sign conventions spelled out in \cite[Section 1.3.4]{hoch2}.




\subsection{Action of the c/o structure of discretely weighted arc-graphs}


We say a c/o structure acts via correlation functions if the
composition of the elements of the c/o is compatible with the
composition of the correlation functions. In short $Y(\alpha
\circ_{w,w'} \alpha')=Y(\alpha)  \circ_{w,w'} Y(\alpha')$ where
$\circ_{w,w'} $ denotes the gluing of the window $w$ and $w'$ holds
as well as the corresponding equation for the self-gluing.


\begin{thm}
\label{discretethm} For a basic  $\B$-Frobenius algebra the c/o
structure of discretely weighted arc-graphs acts  on  the
collection of complexes $B(\beta)$ and the isomorphic Hochschild
complexes via the correlation functions $Y$.
\end{thm}

Just like there are algebras over operads, we can def\/ine algebras
over c/o structures.  The theorem above reads: The collection of bar
complexes $B(S,T)$ form an algebra over the c/o structure of
discretely weighted arc-graphs.

\begin{proof}
The proof is a case by case study which occupies Section~\ref{casessection}.
\end{proof}



\subsection{Case by case analysis  of the discrete arc-graph action}
\label{casessection}

The gluing of the surfaces with discrete arcs breaks down into
individual local gluings of pairs of surfaces $S_i$ and $S'_j$. In
case the surfaces are distinct, there are f\/ive cases of this gluing
depicted in Figs.~\ref{gluings} and~\ref{gluingstwo}. The
cases~a) and~b) are the ones familiar from the closed gluing~\cite{hoch2}, the cases~c) and~d) are new in the open/closed case.
The gluing~c) appears when we are gluing two open windows where none
of them is the only window in its boundary component. The gluing~d)
appears when gluing two open windows each of which is the only
window in its boundary component. The most complicated case is when
one of the windows is the only window, while the other is not. This
case~e) is given in Fig.~\ref{gluingstwo}.


\begin{figure}[th]
%\epsfxsize = \textwidth \epsfbox{gluingsnew.eps}
\centerline{\includegraphics[scale=0.96]{Kaufmann-Fig2}}
\caption{Types of local gluings: a)~two sides
without marked points; b)~side with marked point labelled by $\varnothing$
to side with marked point labelled by $\varnothing$; c)~half a labelled
side with marked point labelled by $S$ to half side with marked
point labelled by $T$; d)~full side with labelled point marked by $S$
to full side with labelled point marked by $T$.}\label{gluings}
\end{figure}



\begin{figure}[th]
%\epsfxsize = \textwidth \epsfbox{gluingstwo.eps}
\centerline{\includegraphics[scale=0.96]{Kaufmann-Fig3}}
\caption{e)~Gluing of a lone window with an open
marked point to window with two marked points.}\label{gluingstwo}
\end{figure}

\subsubsection{Non-self gluing, the closed case}
In the case that there is no self-gluing: a)--b) correspond to:
\begin{equation}
\bigg[\int \dots w  \Delta^{\varnothing}_i  w' \dots \bigg] g^{ij}_{\varnothing}
\bigg[\int  \dots w''  \Delta_j^{\varnothing} w''' \dots\bigg]
=\int w' \dots w'' \dots w''',
\end{equation}
where the $w$, $w'$, $w''$ are the weights of the adjacent sides and we use Einstein summation conventions. We have also included
any factors of~$e$ into the ellipses.


\subsubsection{Non-self gluing case in the simple brane case}

Assuming the simple brane case, all labels have to be the same, say
$S$ then c) corresponds to
\begin{gather*}
\bigg[\int \dots w  \idagger_{S}(a_S \Delta^{S}_i  w'\dots\bigg]  g^{ij}_{S}
\bigg[\int  \dots w'' \idagger_S (\Delta_j^{S}a'_S) w''' \dots\bigg]  \\
\qquad{} = \bigg[\int_S \res _S (w' \dots w) a_S \Delta^{S}_i \bigg] g^{ij}_{S}
\bigg[\int_S  \Delta_j^{S}a'_S \res_S( w''' \dots w'')\bigg]\\
\qquad{} = \int w' \dots w \idagger_S (a_S a'_S)  w''' \dots w''
 =  \int \dots w \idagger_S(a_S a'_S)  w'''  \dots w'' w',
\end{gather*}
while d) corresponds to
\begin{gather*}
\bigg[\int \dots w  \idagger_{S}(\Delta^{S}_i  \Delta^{S}_{k}) w'\dots \bigg] g^{ij}_{S}g^{kl}_{S}
 \bigg[\int  \dots w'' \idagger_{S}( \Delta_l^{S} \Delta_j^{S}) w''' \dots \bigg]  \\
\qquad{} = \int w' \dots w \idagger_S (\Delta_i^S g_S^{ij}\Delta_j^S)  w''' \dots w''
 = \int \dots w \idagger_S(e_S)  w'''  \dots w'' w'.
\end{gather*}

The case e) has two subcases. 1) there are three surfaces which are
glued:
\begin{gather*}
\bigg[\!\int \dots w  \idagger_{S}(a_S \Delta^{S}_{i}) w'\dots \bigg]
g^{ij}_{S}
\bigg[\!\int \dots w''  \idagger_{S}( \Delta^{S}_{k} a'_S) w'''\dots \bigg] g^{kl}_{S}
\bigg[\!\int  \dots w^{\rm (vi)} \idagger_{S}( \Delta_j^{S} \Delta_l^{S}) w^{\rm (v)} \dots \bigg]   \\
\qquad{}\times\bigg[ \int_S a'_S\res_S( w'''\dots w'' )\Delta^{S}_{k} \bigg ] g^{kl}_{S}
\int _S  \Delta_l^S \res_S\big(w^{\rm (v)} \dots w^{\rm (vi)}\big)\Delta_j^{S}
  g^{ij}_{S}\bigg[\int_S  \Delta^{S}_{i}\res_S( w' \dots w)a_S\bigg]   \\
\qquad{} = \int_S \dots w\idagger_S(a_Sa'_S) w'''\dots w'' w^{\rm (v)} \dots w^{\rm (vi)} w',
\end{gather*}
and the case 2) where only two surfaces are glued %better words
\begin{gather*}
\bigg[\int \dots w  \idagger_{S}(a_S \Delta^{S}_{i}) w' \dots w''
\idagger_{S}\big( \Delta^{S}_{k} a'_S\big) w'''\dots \bigg] g^{ij}_{S}
g^{kl}_{S} \bigg[\int  \dots w^{\rm (vi)} \idagger_{S}\big( \Delta_j^{S} \Delta_l^{S}\big) w^{\rm (v)} \dots \bigg]   \\
\qquad{} =\bigg[
 g^{ij}_{S} \int_S a'_S\res_S\big( w'''\dots w  \idagger_{S}\big(a_S \Delta^{S}_{i}\big) w' \dots w'' \big)\Delta^{S}_{k}  \bigg] g^{kl}_{S}
 \int _S \Delta_l^S \res_S\big(w^{\rm (v)} \dots w^{\rm (vi)} \big) \Delta_j^{S}
 \\
\qquad{} = g^{ij}_{S} \int_S a'_S\res_S\big( w'''\dots w  \idagger_{S}\big(a_S \Delta^{S}_{i}\big) w' \dots w'' w^{\rm (v)} \dots w^{\rm (vi)} \big) \Delta_j^{S} \\
\qquad{}=  g^{ij}_{S}\int  \dots w  \idagger_{S}\big(a_S \Delta^{S}_{i}\big) w' \dots w'' w^{\rm (v)} \dots w^{\rm (vi)} \idagger\big(\Delta_j^S a_S'\big)w''' \dots.
\end{gather*}

In this case we have to use commutativity $(C)$ and
the Euler compatibility $(E)$:
\begin{gather*}
  g^{ij}_{S}\int  \dots w  \idagger_{S}\big(a_S \Delta^{S}_{i}\big) w' \dots w'' w^{\rm (v)} \dots w^{\rm (vi)} \idagger_S\big(\Delta_j^S a_S'\big)w''' \dots \\
\qquad {}\times g^{ij}_{S}\int  \dots w   w' \dots w'' w^{\rm (v)} \dots w^{\rm (vi)}\idagger_{S}\big(a_S \Delta^{S}_{i}\big) \idagger_S\big(\Delta_j^S a_S'\big)w''' \dots \\
\qquad{} =\int \dots w  w' \dots w'' w^{\rm (v)} \dots w^{\rm (vi)} \idagger_S(a_Sa'_S) e,
\end{gather*}
which is the contribution we get from the glued surfaces, since there is one self-gluing involved  and
this makes the Euler characteristic go up by one.


\subsubsection{The self gluing cases}
\label{twocasespar} So far we have assumed that the two (half) sides
that are glued are on dif\/ferent $S_i$ and $S'_j$. It can happen that
they belong to the same surface. In these cases, much like in the
case e) 2), there are fewer integrals and instead an Euler class
factor $e$ appears. In the case that there is self-gluing: a)--b)
correspond to:
\begin{equation*}
\bigg[g^{ij}_{\varnothing}  \int \dots w  \Delta^{\varnothing}_i  w'    \dots w''  \Delta_j^{\varnothing} w''' \dots\bigg]
=\int w' \dots w'' \dots w''' e.
\end{equation*}
This r.h.s.\ is the contribution to  correlator for the glued
surface since the self-gluing
 changes the Euler characteristic by $-1$. Again we need to use $(C)$.

 The cases c), d) are analogous to the case e) 2). The self gluing decreases the Euler characteristic by one,
 while the summation gives the factor $e$ by condition~$(E)$. The two cases for e) then either involve only one or two integrals, respectively; correspondingly the gluing then gives rise to a factor of $e$ or $e^2$,
respectively.

There is one more local gluing which comes from gluing consecutive windows corresponding
to Fig.~\ref{stropfig4} cases e) and f). In these cases two half sides of a {\em single}
side marked by a point labelled by some $S\neq \varnothing$ are glued together.


\begin{figure}[th]
%\epsfxsize = \textwidth \epsfbox{gluingsthree.eps}
\centerline{\includegraphics[scale=0.97]{Kaufmann-Fig4}}
\caption{(g) Gluing consecutive half sides.}\label{gluingsthree}
\end{figure}

This corresponds to the following equation in the contribution to the correlator of $S_i$
\begin{equation*}
\int \dots w \idagger\big(\Delta_i^Sg^{ij}\Delta_j^S\big)w' \dots=
\int \dots w \idagger(e_S) w'' \dots.
\end{equation*}

When performing the gluing on the whole window, it might happen,
that there are also self-gluings for the surfaces $S_i$ at some
other sides, in which case we need $(C)$ and $(E)$ and proceed as above.
If there are more consecutive open gluings on one $S_i$, we produce
two punctures and correspondingly two factors of~$\idagger_S(e_S)$.


\setcounter{section}{8}
\setcounter{subsection}{0}
\setcounter{subsubsection}{0}
\setcounter{thm}{0}
\setcounter{equation}{6}
% Original source lines 2031--2099
\appendix

\section{Appendix: operadic, PROPic and c/o structures}\label{appendixA}
% use *-form to suppress numbering



\subsection{The def\/inition of a c/o structure}
\label{coapp}
 Specify an object ${\mathcal O}(S,T)$ in some f\/ixed
symmetric monoidal category for each pair $S$ and $T$ of f\/inite
sets.  A {\it $G$-coloring} on ${\mathcal O}(S,T)$ is the further
specif\/ication of an object $G$ in this category and a morphism
$\mu : S\sqcup T \to {\rm Hom}({\mathcal O}(S,T),$G$)$, and we shall let
${\mathcal O}_\mu (S,T)$ denote this pair of data.

A $G$-colored ``closed/open'' or {\it c/o structure} is a
collection of such objects ${\mathcal O}(S,T)$ for each pair of
f\/inite sets $S$, $T$ together with a choice of weighting $\mu$ for
each object supporting the following four operations which are morphisms in the
category:


\vskip .1in

\noindent
\noindent{\it Closed gluing}:  $\forall\, s\in S$, $\forall\, s'\in S'$ with $\mu (s)=\mu '(s')$,
\[
\circ_{s,s'}: \  {\mathcal O}_\mu (S,T)\otimes{\mathcal O}_{\mu '} (S',T')\to {\mathcal O}_{ \mu ''}
(S\sqcup S'-\{ s,s'\}
, T\sqcup T');
\]


\noindent{\it Closed self-gluing}: $\forall \, s,s'\in S$ with $\mu (s)=\mu (s')$ and $s\neq s'$,
\[
\circ^{s,s'}: \ {\mathcal O}_\mu (S,T)\to {\mathcal O}_{\mu ''} (S-\{ s,s'\} , T);
\]



\noindent{\it Open gluing}:
$\forall  \, t\in T$, $\forall \, t'\in T'$ with $\mu (t)=\mu '(t')$,
\[
\bullet_{t,t'}: \ {\mathcal O}_\mu (S,T)\otimes{\mathcal O}_{\mu '} (S',T')\to {\mathcal O}_{\mu ''}
(S\sqcup S' , T\sqcup T'-\{t,t'\});
\]



\noindent{\it Open self-gluing}: $\forall \, t,t'\in T$ with $\mu (t)=\mu (t')$ and $t\neq t'$,
\[
\bullet^{t,t'}:  \ {\mathcal O}_\mu (S,T)\to {\mathcal O}_{\mu ''} (S , T-\{ t,t'\}).
\]


\noindent In each case, the coloring $\mu ''$ is  induced in the
target in the natural way by restriction, and we assume that $S\sqcup S'\sqcup T\sqcup T'-\{
s,s',t,t'\}\neq\varnothing$.

The  axioms are that the operations are equivariant
for bijections of sets and for bijections of pairs of sets, and the collection of all operations taken together
satisfy associativity.


Notice that we use the formalism of operads indexed by f\/inite
sets rather than by natural numbers as in \cite{MSS} for instance.


\begin{thebibliography}{99}
\bibitem[8]{hoch2}
Kaufmann R.M.,
Moduli space actions on the Hochschild co-chains of a Frobenius algebra. II.~Correlators,
\href{http://dx.doi.org/10.4171/JNCG/22}{{\it J. Noncommut. Geom.}} {\bf 2} (2008), 283--332,
\href{http://arxiv.org/abs/math.AT/0606065}{math.AT/0606065}.


\bibitem[11]{KP}
Kaufmann R.M., Penner R.B.,
Closed/open string diagrammatics,
\href{http://dx.doi.org/10.1016/j.nuclphysb.2006.03.036}{{\it Nuclear Phys.~B}} {\bf 748} (2006), 335--379,
\href{http://arxiv.org/abs/math.GT/0603485}{math.GT/0603485}.

\bibitem[12]{KLP}
Kaufmann R.M., Livernet  M., Penner R.B.,
Arc operads and arc algebras,
\href{http://dx.doi.org/10.2140/gt.2003.7.511}{{\it Geom. Topol.}} {\bf 7} (2003), 511--568,
\href{http://arxiv.org/abs/math.GT/0209132}{math.GT/0209132}.

\bibitem[14]{MSS}
 Markl M., Shnider S., Stashef\/f J.,
Operads in algebra, topology and physics,
{\it Mathematical Surveys and Monographs}, Vol.~96, Amer. Math. Soc., Providence, RI, 2002.


%\bibitem{Merk}
%Merkulov S.A.,
%De Rham model for string topology,
%\href{http://dx.doi.org/10.1155/S1073792804132662}{{\it  Int. Math. Res. Not.}} {\bf  2004} (2004),  no.~55, 2955--2981,
%\href{http://arxiv.org/abs/math.AT/0309038}{math.AT/0309038}. !!!Not Cited in the paper!!!


\bibitem[15]{P1}
Penner R.C.,
Decorated Teichm\"uller theory of bordered surfaces,
\href{http://projecteuclid.org/getRecord?id=euclid.cag/1098468019}{{\it Comm. Anal. Geom.}} {\bf  12}  (2004),  793--820,
\href{http://arxiv.org/abs/math.GT/0210326}{math.GT/0210326}.

\bibitem[16]{Qu}
Quillen D.,
Elementary proofs of some results of cobordism theory using Steenrod operations,
\href{http://dx.doi.org/10.1016/0001-8708(71)90041-7}{{\it Adv. Math.}} {\bf 7} (1971), 29--56.

\end{thebibliography}
\end{document}
